\section{Polynomial Rings and Ideals}

The study of systems of multivariate polynomial equations is a cornerstone of both computational algebra and algebraic geometry. Consider a typical system of equations
\begin{equation}
\label{eq:system1}
\begin{cases}
    \ f_{1}(x,y) = x^{2} - y^{3} = 0\\
    \ f_{2}(x,y) = x^{2} - 4y + 3 = 0
\end{cases}
\end{equation} \\
which arise naturally across many diverse fields, like robotics \cite{usage_robotics1}, signal processing \cite{usage_signalprocessing1}, and statistics \cite{usage_stats1,usage_stats2}. The most fundamental question regarding such a system is the \textit{consistency} problem, which asks whether there exists an exact common solution $(x_0, y_0) \in \mathbb{C}^2$ such that all equations are simultaneously satisfied.

To approach this problem algebraically, we work within a formal structure known as a polynomial ring.

\begin{definition}[Polynomial Ring]
\label{def:poly_ring}
The polynomial ring $R = k[x_1, \dots, x_n]$ is the set of all polynomials in variables $x_1, \dots, x_n$ with coefficients from a field $k$ (typically $\mathbb{Q}, \mathbb{R},$ or $\mathbb{C}$). $R$ is an associative algebra equipped with standard addition and multiplication operations for polynomials.
\end{definition}

One powerful way to show that a system has no solution is to find a contradiction through a linear combination of its generators. If we can find polynomials $a_1, a_2 \in R$ such that:
\begin{equation}
\label{eq:combination}
a_{1}(x,y)f_{1}(x,y) + a_{2}(x,y)f_{2}(x,y) = 1,
\end{equation}
then the system must be inconsistent. Indeed, any hypothetical solution $(x_0, y_0)$ would force the left-hand side of (\ref{eq:combination}) to zero, contradicting the fact that the right-hand side is the constant $1$. According to Hilbert's Nullstellensatz, over an algebraically closed field like $\mathbb{C}$, the converse is also true; as in, a system has no common zeros if and only if $1$ can be expressed as a polynomial combination of the generators.

This observation shifts our focus from the equations themselves to the set of all possible polynomial combinations they can generate. This set is known as an ideal.

\begin{definition}[Ideal]
\label{def:ideal}
Let $f_1, \dots, f_s$ be polynomials in $R$. The ideal generated by these polynomials, denoted $\langle f_1, \dots, f_s \rangle$, is the set
\begin{equation}
\langle f_1, \dots, f_s \rangle = \left\{ \sum_{i=1}^{s} a_i f_i \mid a_i \in R \right\}.
\end{equation}
\end{definition}

Algebraically, an ideal is a sub-module of $R$ that is closed under addition and under multiplication by any element of $R$. Conceptually, an ideal $I$ represents the collection of all polynomial consequences of its generators; if a point is a zero for every generator of $I$, it is necessarily a zero for every polynomial in $I$.

The consistency of the system in (\ref{eq:system1}) is therefore equivalent to the ideal membership problem; determining whether the constant polynomial $1$ lies within the ideal $I = \langle x^{2} - y^{3}, x^{2} - 4y + 3 \rangle$. While it may not be immediately obvious for this specific example, the system (\ref{eq:system1}) does in fact possess solutions (for instance, $(-1, 1)$), meaning $1 \notin I$. We could easily imagine more complex systems of polynomial equations, where assessing membership is far less straightforward. Thus, developing a systematic algorithmic test for such membership is the primary motivation for the theory of Gröbner bases.